\append{Some useful upper bounds}
\label{app:upper-bounds}

We present here an upper bound for the squared absolute value of 
$\alpha(\ell,\hat{\mathbf{u}}) \, \alpha(m,\hat{\mathbf{v}}) \, \alpha(n,\hat{\mathbf{w}})$,
that is defined by equation \ref{eq:diff-operators-generic} using $\ell, m , n \in \{0, 1, 2\}$
and constant unit vectors $\hat{\mathbf{u}}$, $\hat{\mathbf{v}}$ and $\hat{\mathbf{w}}$. 
To do this, we first convert the wavenumber $k_{x}$ and $k_{y}$ to polar coordinates 
$k_{x} = k \, \cos\theta$, $k_{y} = k \, \sin\theta$ and $k = \sqrt{k_{x}^{2} + k_{y}^{2}}$
to show that
\begin{equation}
	\begin{split}
		\vert \alpha(\ell,\hat{\mathbf{u}}) \, \alpha(m,\hat{\mathbf{v}}) \, \alpha(n,\hat{\mathbf{w}}) \vert^{2} 
		&= \vert (\mathbf{k} \cdot \hat{\mathbf{u}})^{\ell} \, (\mathbf{k} \cdot \hat{\mathbf{v}})^{m} \, (\mathbf{k} \cdot \hat{\mathbf{w}})^{n} \vert^{2} \\
		&= \left[ (\mathbf{k} \cdot \hat{\mathbf{u}})^{2} \right]^{\ell} \left[ (\mathbf{k} \cdot \hat{\mathbf{v}})^{2} \right]^{m} 
		\left[ (\mathbf{k} \cdot \hat{\mathbf{w}})^{2} \right]^{n} \\
		&\le \left[ k_{x}^{2} + k_{y}^{2} + k^{2} \right]^{\ell} \left[ k_{x}^{2} + k_{y}^{2} + k^{2} \right]^{m} 
		\left[ k_{x}^{2} + k_{y}^{2} + k^{2} \right]^{n} \\
		&\le 3^{q} k^{2 \, q} \: , \quad q = \ell+m+n \: ,
		\quad \ell, m , n \in \{0, 1, 2\} \: .
	\end{split}
\label{eq:alpha-upper-bound}
\end{equation}

We are also interested in defining an upper bound for the absolute value
of $\mathcal{F} [ U(\mathbf{r}^{\dagger}) ]$ (equation \ref{eq:potential-U-Fourier}), 
i.e., for the amplitude spectrum of $U(\mathbf{r})$ (equation \ref{eq:potential-U-space}). 
To do this, we have adapted the approach presented by
\citet{bhattacharyya1967} to define an upper bound for the amplitude spectrum of 
$\tfrac{\partial}{\partial z} U(\mathbf{r})$. 
We first stress that 
$e^{- k \, (z' - z^{\dagger}_{0})} \le e^{- k \, (z^{t}_{0} - z^{\dagger}_{0})}$ because
$z^{t}_{0} \le z'$, where $z^{t}_{0}$ defines the top of the sources and $z'$ the vertical
coordinate of a point inside the sources. Using this identity into  
$\mathcal{F} [ U(\mathbf{r}^{\dagger}) ]$ (equation \ref{eq:potential-U-Fourier}), we obtain
\begin{equation}
	\vert \mathcal{F} [ U(\mathbf{r}^{\dagger}) ] \vert \le 2 \pi \kappa 
	\frac{e^{- k \, \Delta z}}{k}
	\int\limits_{z^{t}_{0}}^{\infty}
	\mathcal{F} [ \sigma(\mathbf{r}') ] \:
	\: dz' \: , \quad \Delta z = z^{t}_{0} - z^{\dagger}_{0} > 0 \: .
	\label{eq:potential-U-spectrum-bound1}
\end{equation}
Manipulating $\mathcal{F} [ \sigma(\mathbf{r}') ]$ (equation \ref{eq:sigma-Fourier}), 
we obtain
\begin{equation}
	\vert \mathcal{F} [ \sigma(\mathbf{r}') ] \vert \le
	\int\limits_{-\infty}^{\infty}\int\limits_{-\infty}^{\infty}
	\vert \sigma(\mathbf{r}') \vert 
	\: dx'dy'
	\label{eq:sigma-spectrum-bound}
\end{equation}
and, consequently, that
\begin{equation}
	2 \pi \kappa \, 
	\int\limits_{z^{t}_{0}}^{\infty}
	\vert \mathcal{F} [ \sigma(\mathbf{r}') ] \vert
	\: dz' \le \Omega \: ,
	\label{eq:Omega-bound}
\end{equation}
where
\begin{equation}
	\Omega = 2 \pi \kappa \, \iiint\limits_{\mathcal{V}}
	\vert \sigma(\mathbf{r}') \vert 
	\: d\mathcal{V}'
	\label{eq:Omega}
\end{equation}
is a constant depending only on the position and shape of the sources. 
Then, by combining equations 
\ref{eq:potential-U-spectrum-bound1}--\ref{eq:Omega},
we conclude that the amplitude spectrum of 
$U(\mathbf{r})$ (equation \ref{eq:potential-U-space}) on a horizontal plane 
$z^{\dagger}_{0} < z^{t}_{0}$ located above the sources satisfies 
\begin{equation}
	\vert \mathcal{F} [ U(\mathbf{r}^{\dagger}) ] \vert \le
	\frac{e^{- k \, \Delta z}}{k} \, \Omega \: ,
	\: \quad \Delta z = z^{t}_{0} - z^{\dagger}_{0} > 0 \: .
	\label{eq:potential-U-spectrum-bound}
\end{equation}
From the upper bound of $\vert \mathcal{F} [ U(\mathbf{r}^{\dagger}) ] \vert$
(equation \ref{eq:potential-U-spectrum-bound}) and $\mathcal{F} [ d^{o}(\mathbf{r}^{\dagger}) ]$
(equation \ref{eq:observed-data-Fourier}), we also deduce
\begin{equation}
	\vert \mathcal{F} [ d^{o}(\mathbf{r}^{\dagger}) ] \vert \le
	\vert \alpha_{d} \vert \, \Omega \, 
	\frac{e^{- k \, \Delta z}}{k} \: ,
	\: \quad \Delta z = z^{t}_{0} - z^{\dagger}_{0} > 0 \: .
	\label{eq:observed-data-spectrum-bound}
\end{equation}
Note that, for the case in which $\alpha_{d} = k$ (equation \ref{eq:observed-data-Fourier}), 
we obtain the same result presented by \citet{bhattacharyya1967} as the upper bound for the 
amplitude spectrum of $\tfrac{\partial}{\partial z} U(\mathbf{r})$, whereas
$\alpha_{d} = 1$ leads to the upper bound defined by equation \ref{eq:potential-U-spectrum-bound}.
Equations \ref{eq:potential-U-spectrum-bound} and \ref{eq:observed-data-spectrum-bound} 
show that the amplitude spectra 
$\vert \mathcal{F} [ U(\mathbf{r}^{\dagger}) ] \vert$ and 
$\vert \mathcal{F} [ d^{o}(\mathbf{r}^{\dagger}) ] \vert$ on a plane $z^{\dagger}_{0} < z^{t}_{0}$ 
above the sources 
(i) are finite at all wavenumbers $k = \sqrt{k_{x}^{2} + k_{y}^{2}}$ and (ii) tend
to zero as $k$ goes to infinite.
These results confirm the well-known fact that 
$\mathcal{F} [ U(\mathbf{r}^{\dagger}) ]$ and $\mathcal{F} [ d^{o}(\mathbf{r}^{\dagger}) ]$ 
(equations \ref{eq:potential-U-Fourier} and \ref{eq:observed-data-Fourier}) are square integrable.